\chapter{Python for the .NET Engineer}
\label{ch:python}

\begin{quote}\itshape
The language is the small part. The idioms are what a reviewer notices.
\end{quote}

\noindent
You can be productive in Python in a fortnight and still write code that a Python
reviewer will quietly rewrite. This chapter is about the difference. It is not a
tour of the language; it is a list of the specific places where an engineer with
a decade of C\# behind them reliably produces something that works and reads
wrong --- and of the parts of the modern Python toolchain that have no C\#
equivalent and therefore get skipped.

\chapterstub{
\textbf{Argument.} Python's type system is a linting convention rather than a
runtime contract, and Pydantic is what fills the gap. Understanding \emph{that}
one fact explains most of the design of the libraries in this book: schemas,
structured output, state definitions and settings are all Pydantic or
\code{TypedDict} underneath.

\textbf{Sections.} (2.1) Environment: \code{uv}, \code{pyproject.toml}, lockfiles,
virtual environments, and why \code{uv} rather than the four alternatives.
(2.2) The quality gate: \code{ruff}, \code{mypy} in strict mode, \code{pytest} ---
configured once, in \code{pyproject.toml}. (2.3) Typing that matters here: type
hints, \code{Protocol} (structural, unlike a C\# interface), \code{TypedDict},
\code{Annotated}, \code{Literal}, \code{Generic}. (2.4) Pydantic v2: models,
validation, \code{BaseSettings}, serialisation, JSON Schema generation --- and
where each of the four shows up later in the book. (2.5) Idioms: context
managers, decorators, \code{dataclasses}, comprehensions, ABCs, \code{functools}.
(2.6) Project structure, dependency injection without a container, configuration,
structured logging. (2.7) The concept map, .NET to Python --- the full table.
(2.8) The review checklist: what reads as C\#-in-Python, with the fix for each.

\textbf{Listings.} A complete \code{pyproject.toml} for the mini-project;
\code{ruff}/\code{mypy} configuration; a \code{Protocol} used where a C\#
developer would reach for an interface; a Pydantic settings class; six paired
before/after snippets for the review checklist.

\textbf{Diagrams.} \code{py-toolchain} --- what runs when, from \code{uv sync}
to \code{pytest}. \code{py-typing-map} --- where each typing construct is used by
the libraries in this book.

\textbf{Mini-project.} Stage 00: the repository skeleton, tooling configured,
one trivial test passing, CI green.

\textbf{Source topics covered.} Part~I 1.1--1.7 (all of which the source marks
as gaps).
}
